\section*{Chapter \ref{ch:iv1:time-series} --- Time Series}
\markboth{\protect\omTocMark{}{Answers: Time Series}}{\protect\omTocMark{}{Answers: Time Series}}

\begin{solution}{iq:iv1:time-series:1}
Stationary: (a), (c), (e), (f). Not stationary: (b) and (d).

Weak stationarity\index{stationarity} needs a constant mean, a finite constant variance and an autocovariance that depends only on the lag.

(a) holds by definition. (c) holds, as $|\phi| < 1$ and it starts in its stationary law. (e) is white noise, the increments of the walk. (b) has variance growing with $t$. (d) has a changing mean, though removing the trend makes it stationary.

(f) is the surprise. Its mean is 0, its variance is $\cos^2t + \sin^2t = 1$ at every $t$, and its covariance is $\cos(t - s)$, which depends only on the lag. So it is stationary, though every path is a perfect wave. It is not ergodic: one path never reveals that $A$ and $B$ could be different, so the time average of $x_t^2$ gives $(A^2 + B^2)/2$ and not the variance 1, although the time average of $x_t$ does recover the mean 0.

\iqwrong{Calling (f) non-stationary because it has a visible pattern, or (d) stationary because it looks stable about its trend.}

\iqfollow{Why does non-ergodicity matter in practice?}{Time averages from a single history then need not estimate the ensemble values. Only one path of the market is ever observed.}

\iqwhy{Stationarity is often stated and rarely tested at the definition. The wave shows that a pattern and stationarity can coexist.}
\end{solution}

\begin{solution}{iq:iv1:time-series:2}
The variance is $1.96$, the lag-3 autocorrelation $0.343$, the half-life $1.94$ steps, and the $R^2$ is $0.7^{2h}$: $0.49$ at $h = 1$ and $0.24$ at $h = 2$.

Taking the variance of both sides gives $\gamma_0 = 0.49\gamma_0 + 1$, so $\gamma_0 = 1/0.51$\index{AR(1) process}. The autocorrelation at lag $k$ is $\phi^k$. A shock decays as $\phi^h$, so the half-life is $\ln 2/(-\ln 0.7)$.

The best forecast of $x_{t + h}$ is $\phi^hx_t$. It has variance $\phi^{2h}\gamma_0$, so the share of variance explained is $\phi^{2h}$. A long simulation confirms the variance, the autocorrelation and the second $R^2$.

As $\phi \to 1$, the half-life and the variance blow up, and the process approaches a random walk.

\iqwrong{Using $1/(1 - \phi)$ for the variance, which gives $3.3$. It is $1/(1 - \phi^2)$.}

\iqfollow{What is the error variance of the 2-step forecast?}{$1 + \phi^2 = 1.49$, which equals $\gamma_0(1 - \phi^4)$. Long-horizon error approaches the unconditional variance.}

\iqwhy{Quick fluency with the workhorse model. The $R^2$ formula makes clear how fast forecastability fades.}
\end{solution}

\begin{solution}{iq:iv1:time-series:3}
The variance is $1.25$, the lag-1 autocorrelation $0.4$, and the lag-2 autocorrelation 0. The largest possible value is $\frac12$, and the coefficient 2 gives the same autocorrelations.

The variance is $(1 + \theta^2)\sigma^2$. The lag-1 covariance is $\theta\sigma^2$, as only $\varepsilon_{t - 1}$ is shared, so $\rho_1 = \theta/(1 + \theta^2) = 0.5/1.25 = 0.4$\index{MA(1) process}. Beyond lag 1 there is no overlap, so autocorrelations vanish. A simulation agrees.

$\theta/(1 + \theta^2)$ peaks at $\theta = 1$ with value $\frac12$. So any lag-1 autocorrelation over $\frac12$ with nothing beyond it cannot come from an MA(1).

Replacing $\theta$ with $1/\theta$ keeps the ratio. For $\theta = 2$: $2/5 = 0.4$. The convention picks $|\theta| < 1$, the invertible form\index{invertibility}, in which the noise can be recovered from past data.

\iqwrong{Giving $\rho_1 = 0.5$, the coefficient itself. That holds for an AR(1) but not here.}

\iqfollow{Where does a negative MA(1) in returns come from?}{From noise in observed prices. If $p = x + u$, the differences have a $u_t - u_{t - 1}$ term, and this gives a negative lag-1 autocovariance.}

\iqwhy{This is the mirror image of the AR(1) and the basis for identifying models from correlations.}
\end{solution}

\begin{solution}{iq:iv1:time-series:4}
The latest day has weight 6\%. The half-life is $11.2$ days, the effective number of days about 32, and the volatility estimate jumps by a factor of $1.56$.

The weights are $0.06 \times 0.94^k$, summing to 1\index{exponential weighting}. The half-life solves $0.94^k = \frac12$, so $k = 11.2$.

The effective sample size is $1/\sum w_k^2 = (1 + 0.94)/(1 - 0.94) = 32.3$. The mean age of the data is $1/0.06 = 16.7$ days.

The update is $\hat{\sigma}_{\text{new}}^2 = 0.94\hat{\sigma}^2 + 0.06r^2$. With $r = 5\hat{\sigma}$, this is $(0.94 + 1.5)\hat{\sigma}^2 = 2.44\hat{\sigma}^2$. The volatility rises by $\sqrt{2.44} = 1.56$. So one day has a large effect, and it persists: after 11 days, half of the added variance remains.

\iqwrong{Saying that the volatility rises by $5 \times 6\%$, which mixes up volatility and variance.}

\iqfollow{How does this relate to GARCH?}{It is an integrated GARCH(1,1) with $\omega = 0$, $\alpha = 0.06$, $\beta = 0.94$. There is no long-run variance to return to, so the forecast is flat at every horizon.}

\iqwhy{This is the default risk estimate at many desks, and the jump after a large move is what a risk manager has to explain.}
\end{solution}

\begin{solution}{iq:iv1:time-series:5}
It is stationary, with complex roots $0.6 \pm 0.374i$ of modulus $0.707$ and a period of $11.3$ steps. The autocorrelations are $0.8$ and $0.46$.

Stationarity depends on the roots of $z^2 - 1.2z + 0.5 = 0$\index{AR(2) process}, which are $0.6 \pm \sqrt{0.36 - 0.5}$. Their modulus is $\sqrt{0.5} = 0.707 < 1$, so it is stationary.

The argument is $\arctan(0.374/0.6) = 0.558$ radians, so the cycle has period $2\pi/0.558 = 11.3$.

The Yule--Walker equations\index{Yule-Walker equations} give $\rho_1 = \phi_1 + \phi_2\rho_1$, so $\rho_1 = 1.2/1.5 = 0.8$. Then $\rho_2 = 1.2(0.8) - 0.5 = 0.46$. A long simulation gives the same.

A shock has effects 1, $1.2$, $0.94$, $0.53$, \dots. It first grows, then falls, turns negative at step 5, and oscillates, shrinking by $0.707$ per step, which is a damped wave. An AR(1) with a positive coefficient cannot do this. Note that the coefficient $1.2$ exceeds 1 and is harmless, as what matters is the roots.

\iqwrong{Saying that $\phi_1 > 1$ makes it explosive, or checking only $\phi_1 + \phi_2 < 1$, which leaves out the other conditions.}

\iqfollow{Does a cycle in the model mean a tradable cycle?}{Not by itself. The damping is strong, and the wave's phase wanders with new shocks, so the cycle's timing is lost within a few periods.}

\iqwhy{Reading dynamics from the roots is the key skill. Hump-shaped responses appear in macro and in order-flow data.}
\end{solution}

\begin{solution}{iq:iv1:time-series:6}
An AR(1) with $\phi \approx 0.6$: the autocorrelations decay geometrically, and the partial autocorrelations cut off after lag 1. The bands are $\pm 2/\sqrt{500} = \pm 0.089$.

$0.6^2 = 0.36$, $0.6^3 = 0.22$ and $0.6^4 = 0.13$, which fit the first list. The partial autocorrelation\index{partial autocorrelation} at lag $k$ is the last coefficient in an AR($k$) fit, and it is the correlation left once the lags between are controlled for. After lag 1, all are within the bands, so one lag is enough.

An MA(1) has the reverse: autocorrelation cuts off after lag 1, and the partial ones decay, alternating in sign if the coefficient is positive. An ARMA(1,1) has both decaying, with neither cutting off\index{model identification}.

Simulations confirm the signature: a lag-2 regression coefficient near 0 for the AR(1), and near $-0.19$ for an MA(1) with coefficient $0.5$.

The bands are for white noise. Under the fitted AR(1), the sample autocorrelations at higher lags have a wider spread, so the bands understate the uncertainty there.

\iqwrong{Fitting a high order because the autocorrelations stay above the band for several lags. A slow decay alone says nothing about the order, and the partial values are what show it.}

\iqfollow{How many lags will look significant by chance among 20, with no true structure?}{About one at the 5\% level. Testing many lags together is the role of the Ljung--Box statistic.}

\iqwhy{This is the classic identification table, and it can be asked of real output in a minute.}
\end{solution}

\begin{solution}{iq:iv1:time-series:7}
$1.1$ and $1.173$, with a limit of $(1 + \rho)/(1 - \rho) = 1.222$. True long-horizon volatility is about $10.6\%$ above the square-root-of-time figure.

The variance of a sum of $q$ terms gives\index{variance ratio}
\[
\mathrm{VR}(q) = 1 + 2\sum_{k = 1}^{q - 1}\left(1 - \frac{k}{q}\right)\rho_k.
\]
With $\rho_k = 0.1^k$: $\mathrm{VR}(2) = 1 + 0.1$, and $\mathrm{VR}(5) = 1 + 2(0.8 \times 0.1 + 0.6 \times 0.01 + 0.4 \times 0.001 + 0.2 \times 0.0001) = 1.173$. A long simulation agrees.

For large $q$, the sum tends to $1 + 2\rho/(1 - \rho) = 1.222$. Volatility scales by the root of the ratio, $\sqrt{1.222} = 1.106$. So a small daily autocorrelation adds 10.6\% to annual volatility, and it would turn a measured annual Sharpe ratio of 1 into about $0.9$.

With $\rho < 0$, the ratio drops below 1, as for mean-reverting returns, which look worse on daily data than on long horizons.

\iqwrong{Using $\mathrm{VR}(q) = 1 + (q - 1)\rho$, which ignores both the weights and the decay of $\rho_k$.}

\iqfollow{A fund's monthly returns have autocorrelation $0.3$. What does this suggest?}{Stale marks or smoothing, more likely than skill, and it means that the true volatility is higher by a factor of about $\sqrt{1.857} = 1.36$.}

\iqwhy{This is the simplest test of the random walk, and it has a direct effect on any annualised number.}
\end{solution}

\begin{solution}{iq:iv1:time-series:8}
The long-run variance is $4 \times 10^{-5}$, or 10\% annual volatility, the half-life is $13.5$ days, and the forecast is $2.8$ times the long-run variance, or $16.8\%$.

Taking expectations of the recursion\index{GARCH}, $V = \omega + (\alpha + \beta)V$, so $V = \omega/(1 - \alpha - \beta) = 2 \times 10^{-6}/0.05$. Annualising gives $\sqrt{252V} = 10\%$.

The conditional forecast obeys $\mathbb{E}\sigma_{t + h}^2 - V = (\alpha + \beta)^h(\sigma_t^2 - V)$. So shocks decay at $0.95$ per day, with half-life $\ln 0.5/\ln 0.95 = 13.5$.

Today's excess is $3V$. After 10 days, it is $0.95^{10} \times 3V = 1.8V$, so the forecast is $2.8V$, an annual volatility of $16.8\%$. A Monte Carlo over the recursion gives $2.8$.

If $\alpha + \beta = 1$, $V$ does not exist, and the forecast stays flat.

\iqwrong{Applying the decay to volatility, not variance. The recursion is linear in variance only.}

\iqfollow{What does this imply for the term structure of forecast volatility?}{It rises or falls smoothly toward 10\% as the horizon grows. The average variance over $h$ days is what matters for a term volatility.}

\iqwhy{GARCH(1,1) is the baseline for volatility forecasting. A candidate should read persistence and the long-run level straight off the coefficients.}
\end{solution}

\begin{solution}{iq:iv1:time-series:9}
The daily volatility is $1.19\%$. It is about five times as efficient as one squared return. Discrete sampling biases it low, and bid-ask bounce biases it high.

$\ln 1.02 = 0.0198$, so the estimate is $0.0198^2/(4\ln 2) = 1.41 \times 10^{-4}$, whose root is $1.19\%$\index{Parkinson estimator}.

The constant comes from $\mathbb{E}R^2 = 4\ln 2\,\sigma^2$ for the range $R$ of a driftless Brownian motion. A squared return has relative variance 2. The Parkinson estimator has relative variance about $0.41$, which gives a ratio of $4.9$. In a simulation with fine steps, the ratio is between $4.3$ and $5.5$.

On real data, it goes wrong in two main ways. First, the true high and low fall between prints, so the observed range is too small. With 78 prices a day, a simulation shows the variance estimate about 14\% too low, about 7\% in volatility. Second, prints at the bid and at the ask extend the range, and that pushes it up, more so in illiquid names. Overnight moves are also missed entirely, so the estimator describes trading hours only, and drift inflates it slightly.

\iqwrong{Treating the fivefold gain as free. It relies on the model, and the two biases work against each other in ways that depend on the stock.}

\iqfollow{How would you combine the range and the returns?}{Use estimators of the Garman--Klass type that use open, high, low and close, and add a close-to-open term for the night.}

\iqwhy{Range estimators are cheap and useful, and the interviewer wants the efficiency figure together with the reasons not to trust it.}
\end{solution}

\begin{solution}{iq:iv1:time-series:10}
The mean is $\sigma^2 + 2n\eta^2$. The relative bias is $+39\%$ at 78 returns and $+195\%$ at 390. The noise gives returns a lag-1 autocovariance of $-\eta^2$.

An observed return is $\Delta x + \Delta u$\index{microstructure noise}, with $\Delta u = u_t - u_{t - 1}$ of variance $2\eta^2$. Summing $n$ squared returns gives $\sigma^2$ from the efficient price and $2n\eta^2$ from the noise. The noise term grows with $n$, while the signal does not.

With $\sigma^2 = 10^{-4}$ and $\eta^2 = 2.5 \times 10^{-7}$: $2 \times 78 \times 2.5 \times 10^{-7} = 3.9 \times 10^{-5}$, or $39\%$ of $\sigma^2$. At 390, it is five times as large. A simulation with these numbers matches the formula within 1\%.

Consecutive returns share $u_t$ with opposite signs, so the autocovariance is $-\eta^2$ whatever the frequency. A plot of realised variance against the sampling frequency, the signature plot\index{signature plot}, rises as the steps shrink.

Roll's estimate reads the same effect as a spread: with bounce of half-spread $s/2$, $\eta = s/2$, and $s = 2\sqrt{-\operatorname{Cov}}$. Here this is a spread of $0.1\%$.

\iqwrong{Assuming that more data always means a better estimate. In this setting, the highest frequency gives the worst estimate of all if nothing is done about the noise.}

\iqfollow{How can autocovariance be used to correct realised variance?}{Add twice the first-order autocovariance sum to the realised variance, which removes the $2n\eta^2$ bias in this model. This is the idea behind realised kernels.}

\iqwhy{This is the first thing that goes wrong with high-frequency data. The sign of the autocovariance is a useful diagnostic.}
\end{solution}

\begin{solution}{iq:iv1:time-series:11}
Under a unit root, the regressor is non-stationary, so the statistic has a non-normal limit. The critical values are $-1.95$ and $-2.87$. Expect $\hat{\rho} \approx 0.979$.

The usual asymptotics rely on $\frac{1}{T}\sum x_{t - 1}^2$ settling to a constant. Under a random walk, it grows like $T$ and stays random. So the limit is a ratio of functionals of Brownian motion\index{Dickey-Fuller test}, and it is skewed to the left. A simulation of $40\,000$ paths with $T = 250$ gives the 5\% quantiles as $-1.95$ with no constant and $-2.87$ with a constant.

The estimate itself is biased downward\index{unit root}. With a constant, $\mathbb{E}\hat{\rho} \approx 1 - 5.3/T$, or $0.979$ at $T = 250$, which a simulation confirms. So a random walk looks mean-reverting, with an apparent half-life of about 33 days. A test using the normal table at $-1.645$ would reject the true unit root about 45\% of the time with a constant.

Removing the constant does not solve this: the 5\% value is still far from normal, and it requires a mean of zero, which is rarely a safe assumption.

\iqwrong{Comparing the statistic with $-1.645$. It is the commonest way to discover pairs that do not exist.}

\iqfollow{Why is the bias so much worse with a constant?}{The demeaned series borrows from the future, so the regressor is correlated with later noise. This is a shift of about $5.3/T$ rather than about $1.8/T$.}

\iqwhy{Every mean-reversion trade rests on this test. Knowing the bias and the critical values prevents the commonest error in pairs research.}
\end{solution}

\begin{solution}{iq:iv1:time-series:12}
The statistic is $2.0$ under independence, and $1.63$ with the AR(1). Selection of the best of many models makes even $2.0$ weak.

Diebold--Mariano\index{Diebold-Mariano test} is a $t$-statistic on the mean difference in loss. With independence, it is $0.02/(0.2/\sqrt{400}) = 2.0$, a two-sided $p$-value of $0.046$.

With autocorrelation, the variance of the mean is multiplied by the long-run factor $(1 + \rho)/(1 - \rho) = 1.5$. So the statistic is $2.0/\sqrt{1.5} = 1.63$, a $p$-value of $0.10$. Simulated AR(1) series with these parameters give the same standard error for the mean. Forecast errors at horizons above one step overlap, and they normally have this kind of dependence. So a heteroskedasticity and autocorrelation consistent estimate, with lags to match, is needed.

If one model is the best of 20 candidates, then the expected best null statistic is already about $1.87$. A value of $2.0$ is then hardly evidence, and a reality check of the data-snooping type, or a separate holdout, is needed\index{data snooping}. The test also compares forecasts as given, so if the models are nested, or the estimation noise matters, the limit is not standard.

\iqwrong{Using the plain standard error for overlapping forecasts. A $p$-value that crosses 5\% often no longer does after that correction.}

\iqfollow{What is the difference between in-sample and out-of-sample fit here?}{This test scores out-of-sample forecasts. In-sample fit always favours the larger model, and it would say nothing here.}

\iqwhy{Claims of improvement need a standard error, and comparing forecasts is how models are picked. The two adjustments cover most real cases.}
\end{solution}

\begin{solution}{iq:iv1:time-series:13}
$x$ and $y$ are I(1), and $u$ is I(0). The half-life is $6.6$ steps, and the slope converges at rate $1/T$, not $1/\sqrt{T}$. Engle--Granger regresses, and then tests the residuals with a special table.

$y$ and $x$ wander, while $y - 2x$ does not. That is cointegration\index{cointegration} with vector $(1, -2)$. The half-life of $u$ is $\ln 2/(-\ln 0.9) = 6.6$.

The slope is superconsistent. Since $\sum x^2$ grows like $T^2$, the error $\sum xu/\sum x^2$ shrinks like $1/T$. In a run, quadrupling $T$ reduces the slope's error by a factor of about 4, against 2 in a stationary regression.

Step one: regress $y$ on $x$ and take the residuals. Step two: test the residuals for a unit root. The usual table fails because $\hat{\beta}$ is chosen to make the residuals as small as possible, so they look stationary even when they are not\index{Engle-Granger method}. On two independent walks with $T = 500$, the 5\% critical value comes out near $-3.37$, far from $-1.95$ or $-1.65$.

With two independent walks, the regression is spurious: the slope stays random, and the $t$-statistic rejects zero over 70\% of the time at the 5\% level. Cointegration is the case in which such a regression has a real meaning, and the residual test tells them apart.

\iqwrong{Testing whether the correlation of the levels is high, or testing with the plain Dickey--Fuller table. Correlation of returns or of levels says nothing about whether a spread comes back.}

\iqfollow{What can you do with several assets?}{Use the Johansen method, a system approach that estimates how many stationary combinations exist, up to the number of assets less one.}

\iqwhy{It is the theoretical base of pairs trading. A candidate must know how it differs from correlation and from spurious regression.}
\end{solution}

\begin{solution}{iq:iv1:time-series:14}
The kurtosis is $3.77$, finite when $(\alpha + \beta)^2 + 2\alpha^2 < 1$. The lag-1 autocorrelation of the squares is $0.179$, decaying at $0.95$ per lag.

Let $V = \mathbb{E}\sigma^2$, $s = \alpha + \beta$ and $m = \mathbb{E}\sigma^4$. Square the recursion and use $\mathbb{E}r^4 = 3\mathbb{E}\sigma^4$ and $\mathbb{E}[r^2\sigma^2] = m$, which follow since $r = \sigma z$ with $z$ independent of $\sigma$\index{GARCH}:
\[
m = \omega^2 + 2\omega sV + (3\alpha^2 + 2\alpha\beta + \beta^2)m.
\]
Note that $3\alpha^2 + 2\alpha\beta + \beta^2 = s^2 + 2\alpha^2$. With $\omega = V(1 - s)$, the constant term is $V^2(1 - s)(1 + s)$. So
\[
m = \frac{V^2(1 - s^2)}{1 - s^2 - 2\alpha^2}, \qquad \kappa = \frac{3m}{V^2} = \frac{3(1 - s^2)}{1 - s^2 - 2\alpha^2}.
\]
With $s = 0.95$ and $\alpha = 0.1$: $3 \times 0.0975/0.0775 = 3.77$. The condition for finiteness is a positive denominator. Here $0.9025 + 0.02 < 1$, so it holds. At $\alpha = 0$, $\kappa = 3$, as it must be for a normal.

The squares follow an ARMA(1,1) with autoregressive coefficient $s$, and the given formula has value $0.1 \times 0.1925/0.1075 = 0.179$. Each further lag multiplies this by $0.95$. A simulation of two million steps gives these values, and a return autocorrelation near 0.

So the returns have no linear memory, while their sizes do\index{volatility clustering}. The value $0.179$ looks modest, yet the slow decay gives substantial total persistence.

\iqwrong{Assuming that fat tails need fat-tailed innovations. Here normal innovations give excess kurtosis through the mixing of variances, although real data usually need more.}

\iqfollow{What is the kurtosis if $\alpha$ is doubled with $\alpha + \beta$ unchanged?}{$3 \times 0.0975/(0.0975 - 0.08) = 16.7$. It is close to the edge: at $\alpha \approx 0.221$ it becomes infinite. So the fourth moment is very sensitive to $\alpha$, and sample kurtosis is an unstable target.}

\iqwhy{This derivation checks that the candidate can work with moments of a recursion. It explains two stylised facts at once.}
\end{solution}

\begin{solution}{iq:iv1:time-series:15}
The $R^2$ is $(\kappa/3 - 1)/(\kappa - 1) = 0.093$, even for a perfect forecast. So the low $R^2$ against squared returns says little, and a better proxy for the variance is the fix.

Write $r^2 = \sigma^2z^2$, with $\mathbb{E}z^2 = 1$ and $z$ independent of $\sigma^2$. Then the regression of $r^2$ on $\sigma^2$ has slope 1, and $R^2 = \operatorname{Var}\sigma^2/\operatorname{Var}r^2$.

Using $\mathbb{E}r^4 = 3\mathbb{E}\sigma^4$ and $\mathbb{E}r^4 = \kappa V^2$: $\operatorname{Var}\sigma^2 = (\kappa/3 - 1)V^2$ and $\operatorname{Var}r^2 = (\kappa - 1)V^2$. At $\kappa = 3.77$: $0.257/2.77 = 0.093$. A simulation of the GARCH itself gives this figure when the forecast is the true $\sigma^2$.

So $R^2$ values of 5\% to 10\% are what perfect models would produce, and the critics who cite them are measuring the proxy's noise, which is enormous for a one-observation estimate\index{volatility proxy}. An imperfect model does worse, which is why the formula is an upper bound for this model.

The fix is a less noisy proxy: realised variance from intraday returns. With 78 normal sub-returns, the noise $\chi^2_{78}/78$ has variance $2/78$, and the formula becomes
\[
R^2 = \frac{\kappa/3 - 1}{\kappa/3 - 1 + (\kappa/3)(2/78)} = 0.89.
\]
A simulation agrees with this within $0.05$. Also, compare models with a loss that is robust to proxy noise, either squared error or the quasi-likelihood loss, and avoid absolute errors of volatility, which can favour the wrong model\index{QLIKE loss}.

\iqwrong{Concluding that GARCH does not forecast, from an $R^2$ of a few percent against squared returns. The test had almost no power in the first place.}

\iqfollow{What does $\kappa = 3$ give, and why?}{$R^2 = 0$: variance is constant, so there is nothing to forecast. The formula gives the sensible limit.}

\iqwhy{This is a well-known confusion in volatility research, with a two-line resolution. It shows whether the candidate separates a model's quality from its measurement.}
\end{solution}

\begin{solution}{iq:iv1:time-series:16}
$\rho_k = R^2\phi^k$, so $\rho_1 = 0.009$ and $\rho_{10} = 0.0035$. The return follows an ARMA(1,1) with $\phi = 0.9$ and a moving-average coefficient of about $0.89$. The autocorrelations are too small to see, and yet they add up.

Since the noise is white and independent\index{persistent expected returns}, $\operatorname{Cov}(r_t, r_{t - k}) = \operatorname{Cov}(\mu_t, \mu_{t - k}) = \phi^k\operatorname{Var}\mu$. Dividing by the variance of $r$ gives $R^2\phi^k$.

Geometric decay from lag 1 on is the signature of an ARMA(1,1), and an AR(1) plus white noise is exactly that. In the form $r_t - \phi r_{t - 1} = e_t - \theta e_{t - 1}$, matching the lag-1 autocorrelation gives $\theta \approx 0.891$, which nearly cancels the root at $0.9$. So it looks like white noise.

To detect $\rho_1 = 0.009$ at $t = 2$ takes $(2/0.009)^2 \approx 49\,000$ observations, or about 196 years of daily data. A simulation of four million steps does show it.

The effect is nonetheless real. The sum of the autocorrelations gives a long-run variance ratio of $1 + 2R^2\phi/(1 - \phi) = 1.18$. So long-horizon returns are more predictable than short-horizon ones: the correlations are small but they all have the same sign.

The key is how the predictability is reached. If $\mu_t$ can be observed through a separate signal, the $R^2$ is 1\%. The best forecast from past returns alone reaches only about $0.04\%$, as both the ARMA innovation variance and a steady-state Kalman filter show. So return history is a poor way to find such a signal.

\iqwrong{Concluding from insignificant autocorrelations that returns are unpredictable. Low autocorrelation fits both no signal and a persistent signal under much noise.}

\iqfollow{How do you exploit this with no outside signal?}{Mostly you cannot, at this $R^2$. Long-horizon regressions gain a little, but they mostly reuse the same data and call for overlap corrections.}

\iqwhy{This resolves an apparent paradox in empirical finance, and it shows why external signals matter more than return history.}
\end{solution}

\enlargethispage{-3\baselineskip}%
\begin{solution}{iq:iv1:time-series:17}
The stationary shares are $\frac23$ and $\frac13$, the volatility is $19.1\%$, the kurtosis is $6.17$, and the autocorrelation of the squares decays as $0.85^k$.

For the shares, balance the flows\index{regime switching}: $\frac23 \times 0.05 = \frac13 \times 0.1$.

The variance is a mixture: $\frac23(0.01) + \frac13(0.09) = 0.0367$, so the volatility is $19.1\%$. A calm-state share of two thirds still leaves most variance in the stressed state.

For kurtosis, the normal has $\mathbb{E}r^4 = 3\sigma^4$ in each state:
\[
\kappa = \frac{3(\frac23 \times 0.0001 + \frac13 \times 0.0081)}{0.0367^2} = 6.17.
\]
This much kurtosis comes from mixing alone, with no fat tails in either state.

The chain's memory is its second eigenvalue, $\lambda = 0.95 + 0.9 - 1 = 0.85$\index{Markov chain eigenvalue}. Conditional on the state, $r_t^2$ and $r_{t + k}^2$ are independent, so their covariance is the covariance of the state variances, which decays as $\lambda^k$. The half-life is $4.3$ days. The level is set by the share of $\operatorname{Var}r^2$ due to the state: $\rho_k = 0.204\lambda^k$.

A four-million-step simulation returns $19.1\%$, a kurtosis near $6.2$, and the $0.85$ decay. As the states approach each other, the kurtosis goes to 3, and the autocorrelation to 0.

\iqwrong{Averaging volatilities, $\frac23(10) + \frac13(30) = 16.7\%$. Variances mix, and volatilities do not.}

\iqfollow{How does this model differ from GARCH with the same persistence?}{The memory is that of an unobserved state, not of past returns as such. A large move raises the forecast only through what it says about the state, which can be a sharp jump, capped at the stressed variance, and not an increase proportional to $r^2$. From a prior of $0.05$, a return of 4 calm-state deviations lifts the stressed probability to $0.956$, and the next day's variance forecast by a factor of $4.5$. Filtering the state is an inference problem.}

\iqwhy{Regime models are often mentioned without numbers. These numbers come from a two-state chain and a line of algebra.}
\end{solution}

\begin{solution}{iq:iv1:time-series:18}
$\mathbb{E}\hat{\phi} \approx 0.9115$, so the reported half-life is about $7.5$ steps, against a true $13.5$. More than half the bias comes from estimating the mean, and a jackknife removes most of it.

The bias is $-(1 + 3 \times 0.95)/100 = -0.0385$\index{AR(1) bias}, so the expected estimate is $0.9115$. Half-lives are $\ln 0.5/\ln\phi$: $13.5$ at $0.95$, and $7.5$ at $0.9115$. A simulation of $200\,000$ samples gives a mean estimate of $0.906$ and a median half-life of about $7.9$, so at this length the true bias is a little larger than the formula's.

There are two sources. Even with a known mean, $\hat{\phi}$ is biased by about $-2\phi/T$, because the shock $\varepsilon_t$ enters later regressors. Demeaning adds about $-(1 + \phi)/T$: the mean of the sample includes the later noise terms, so the demeaned $x_{t - 1}$ moves against $\varepsilon_t$, and the estimate falls. Here the formulas give $-0.019$ and $-0.0195$, and the simulation $-0.017$ and $-0.027$. Near 1 the correction is the largest, and with a small $T$ it can make any random walk look like a reverting one.

So a backtest that selects spreads with short half-lives favours exactly those with the most bias, and expected profits are overstated.

To correct it, a half-sample jackknife\index{jackknife} uses $2\hat{\phi} - \frac12(\hat{\phi}_1 + \hat{\phi}_2)$, where the last two are from the two halves. Since the bias is of order $1/T$, this cancels most of it. In the simulation, the mean returns to $0.953$. Direct use of the formula, solving $\hat{\phi} = \phi - (1 + 3\phi)/T$ for $\phi$, gives $0.945$. Simulation-based or median-unbiased estimates are also used. All of these have a larger variance, which is the cost.

\iqwrong{Reporting the half-life from the raw estimate. A half-life is a highly nonlinear function near $\phi = 1$, so small biases have large effects.}

\iqfollow{Solve the formula for this sample.}{$\phi = (\hat{\phi}T + 1)/(T - 3) = (0.9115 \times 100 + 1)/97 = 0.95$, which recovers the truth by construction when the estimate is at the formula's mean.}

\iqwhy{Half-life estimates decide how much capital a pairs book carries. The bias works in the direction that flatters every candidate.}
\end{solution}

\begin{solution}{iq:iv1:time-series:19}
$T(\hat{\rho} - 1) \to \frac12(W_1^2 - 1)/\int_0^1W^2$, and the $t$-statistic tends to $\frac12(W_1^2 - 1)/\sqrt{\int_0^1W^2}$. The numerator is negative with chance $0.683$. The 5\% value is $-1.95$, and the mean is about $-1.78$.

Start with $\hat{\rho} - 1 = \sum x_{t - 1}\varepsilon_t/\sum x_{t - 1}^2$, where $x_t = \sum_{i \le t}\varepsilon_i$.

Squaring, $x_t^2 = x_{t - 1}^2 + 2x_{t - 1}\varepsilon_t + \varepsilon_t^2$. Summing,
\[
\sum_tx_{t - 1}\varepsilon_t = \frac12\left(x_T^2 - \sum_t\varepsilon_t^2\right).
\]
By Donsker's theorem\index{functional central limit theorem}, $x_{\lfloor Ts\rfloor}/(\sigma\sqrt{T}) \to W_s$. And $\sum\varepsilon_t^2/T \to \sigma^2$. So the numerator divided by $\sigma^2T$ tends to $\frac12(W_1^2 - 1)$. The denominator divided by $\sigma^2T^2$ tends to $\int_0^1W^2$.

Taking the ratio, the factors of $T$ give $T(\hat{\rho} - 1)$, and the $\sigma$ cancels. For the $t$-statistic, the standard error is $s/\sqrt{\sum x_{t - 1}^2}$, with $s \to \sigma$, so it tends to the numerator over $\sqrt{\int W^2}$\index{unit root}.

The numerator is $\frac12(\chi_1^2 - 1)$, a shifted chi-square, and it is negative when $|W_1| < 1$, which has probability $0.683$. So $\hat{\rho} < 1$ about two times in three, even with a true unit root. That is the source of the left skew.

A simulation on $T = 500$ and a direct evaluation of the functional on a fine grid agree: the mean of $T(\hat{\rho} - 1)$ is about $-1.78$, and the $t$ quantile is $-1.95$. The normal value $-1.645$ would reject $9.5\%$ of the time, not 5\%. The denominator has mean $\frac12$, and it is random in the limit, which is the whole reason that normality fails.

\iqwrong{Applying a central limit theorem to the numerator alone. The sum is not normal: it is a chi-square shifted. Each shock is independent of its regressor, but the sizes of the terms are set by $x_{t - 1}$, which grows and stays random, so the sum's conditional variance stays random even after scaling.}

\iqfollow{What does a constant change in the limit?}{Write the limit as $\int_0^1W\,dW/\int_0^1W^2$. Then $W$ is replaced by the demeaned $W^\mu = W - \int_0^1W$ in both integrals, so the numerator becomes $\frac12(W_1^2 - 1) - W_1\int_0^1W$. The law shifts left, and the 5\% value moves to about $-2.86$.}

\iqwhy{This is a compact use of Brownian functionals in econometrics, and it is the proof behind the most used test in the field.}
\end{solution}

\begin{solution}{iq:iv1:time-series:20}
The error is about $2\sigma^4/n + 4n^2\eta^4$, with the best $n = (\sigma^4/(4\eta^4))^{1/3} \approx 34$. Two-scale averaging uses all the ticks and has no first-order bias.

With $n$ returns, the bias is $2n\eta^2$, so its square is $4n^2\eta^4$\index{optimal sampling frequency}. The variance comes mainly from the efficient price, and for a constant volatility it is $2\sigma^4/n$. There are also smaller noise terms, of order $n\eta^4$ and $\sigma^2\eta^2$. Setting the derivative to zero gives $n^3 = \sigma^4/(4\eta^4) = 40\,000$, so $n = 34$: about one return per 11 minutes of a 6.5-hour day. A search over the full curve, including the smaller terms, also lands near there.

The sparse estimator discards most of the data. Two-scale realised variance\index{two-scale realised variance} keeps it. Divide $N$ ticks into $k$ offset subgrids, each of about $\bar{n} = (N - k + 1)/k$ returns. Let $\overline{\mathrm{RV}}$ be the average of their realised variances, whose mean is $\sigma^2 + 2\bar{n}\eta^2$, and let $\mathrm{RV}^{\text{all}}$ use all the ticks, with mean about $2N\eta^2$. Then
\[
\mathrm{TSRV} = \overline{\mathrm{RV}} - \frac{\bar{n}}{N}\mathrm{RV}^{\text{all}}
\]
has mean $\sigma^2$, since the all-ticks estimate is essentially an estimate of the noise.

The simulation has $23\,400$ ticks with these parameters, and $k = 300$. The all-ticks estimate has mean about 118 times the truth. TSRV has a mean within 3\% of it. Its root mean squared error is between 10\% and 30\% of $\sigma^2$, and it is below that of the sparse estimator at $n = 34$, as a good estimator using more data should be.

These assume independent noise. Real noise is autocorrelated and depends on order flow, which is one reason for using more elaborate kernels. And trade time often does better than clock time.

\iqwrong{Choosing a 5-minute grid by habit with no reference to the noise. That is often reasonable but it is the result of this calculation for a certain noise level, not a law.}

\iqfollow{What does the best $n$ become if the noise doubles?}{The formula shows $n \propto \eta^{-4/3}$, so $n$ falls by a factor of $2^{4/3} = 2.5$, to about 14 returns a day.}

\iqwhy{A clean trade-off followed by the idea of using the noisy estimator to correct itself. It connects model-based thinking directly to what is done with tick data.}
\end{solution}
